\section{Introduction}

\subsection{The Structural Compression of IT Services}

The IT services industry is undergoing a structural compression that is eliminating the traditional competitive positions available to specialist firms. This compression operates through four simultaneous forces, each commoditizing a different layer of the value chain.

Hyperscalers --- Amazon Web Services, Microsoft Azure, and Google Cloud --- have moved beyond infrastructure provision into managed AI services, industry-specific solutions, and orchestration capabilities \citep{Jacobides2022}. Each layer they add absorbs work that specialist firms previously differentiated on. Large systems integrators --- Accenture, TCS, Capgemini --- commoditize delivery through scale, acquiring dozens of specialist firms per year and building proprietary platforms that standardize what boutique firms once offered as craft \citep{Gawer2022DigitalAge}. Software-as-a-Service vendors --- Oracle, SAP, ServiceNow --- close the application gap with autonomous, AI-native products where clients subscribe directly, bypassing the integrator entirely and converting what were custom implementation projects into configuration exercises.

Most recently, generative AI --- ChatGPT, Claude, Copilot --- commoditizes knowledge work itself: the highest-margin activity in IT services. Strategy advice, architecture reviews, security assessments, and code generation --- activities that specialist firms charged premium rates for because they required human expertise --- are becoming available at near-zero marginal cost. The knowledge asymmetry that justified the consulting premium is collapsing. Software development itself is being compressed: agentic coding tools mean fewer engineers deliver equivalent output, and every client will ask why they need a fifty-person team when ten people with AI agents produce the same results.

For firms operating without platform ownership --- those who cannot compete on infrastructure scale, talent volume, product IP, or generic knowledge --- the strategic literature offers a stark diagnosis. \citet{Porter1985} identifies this as ``stuck in the middle'': firms that achieve neither cost leadership nor differentiation face systematic underperformance. Transaction cost economics \citep{Williamson1991} predicts that such firms will behave opportunistically to survive, maximizing extraction from each client relationship. Ecosystem strategy \citep{Gawer2014, Jacobides2018} theorizes only two viable roles: platform owner or complementor. The non-platform orchestrator is not a recognized category.

Yet empirically, some firms do the opposite of what theory predicts. They invest in building client capabilities beyond contractual scope. They coordinate competing vendors in the client's interest. They resolve organizational paradoxes that the client cannot resolve internally. They orchestrate complex multi-vendor ecosystems without owning the platform. They act as stewards --- absorbing risk and building client independence rather than dependency.

\subsection{The Theoretical Gap}

This observed behavior violates five core theoretical assumptions simultaneously.

First, the opportunism assumption in transaction cost economics \citep{Williamson1991, Williamson1993OpportunismCritics}. TCE assumes suppliers minimize effort and maximize extraction within contractual boundaries. The observed behavior --- investing in client capabilities at the supplier's short-term cost --- contradicts this foundational behavioral assumption. This is not a parameter adjustment; relaxing the opportunism assumption fundamentally changes the theory.

Second, the internal control assumption in dynamic capabilities theory \citep{Teece1997, Teece2018}. Dynamic capabilities are theorized as internal organizational processes for sensing, seizing, and reconfiguring resources. No framework exists for capabilities that operate \emph{across} organizational boundaries on behalf of others --- what we term meta-dynamic capabilities. The supplier builds capabilities inside the client organization, using knowledge accumulated across its portfolio, through a mechanism that existing DC theory cannot describe.

Third, the platform orchestration assumption in ecosystem strategy \citep{Gawer2014, Adner2017, Jacobides2018}. Ecosystem theory recognizes platform owners who orchestrate through architectural control and complementors who participate. The non-platform orchestrator --- a firm that coordinates ecosystem participants through relational governance and contextual knowledge rather than platform ownership --- is not a theoretical category.

Fourth, the intra-firm boundary of stewardship theory \citep{Davis1997TowardManagement, Donaldson1991}. Stewardship theory extends agency theory by explaining why managers act as stewards of organizational interests rather than self-interested agents. This extension has not been applied to inter-organizational relationships where the steward is an external supplier acting in the client's interest without hierarchical authority or equity ownership.

Fifth, the internal resolution assumption in paradox theory \citep{Smith2011}. Organizational paradoxes --- tensions between innovation and stability, exploration and exploitation, change and continuity --- are theorized as resolvable through internal managerial capabilities. No framework exists for external resolution agents who navigate these tensions from outside, enabled by political neutrality and cross-client pattern recognition.

No single existing construct captures the combination of stewardship motive, external capability building, non-platform orchestration, cross-boundary dynamic capabilities, and external paradox resolution. Adjacent constructs address fragments: the relational view \citep{DyerSingh1998} explains inter-firm competitive advantage through relationships but assumes symmetry; boundary spanning captures cross-boundary activity but not the stewardship motive; coopetition \citep{nalebuff1996coopetition} addresses simultaneous competition and cooperation but not the asymmetric steward role.

\subsection{Contextual Intimacy as the Defensible Position}

If the four compression forces eliminate infrastructure, delivery, applications, and generic knowledge as sources of competitive advantage, what remains? We argue that the only defensible position for the non-platform supplier is \emph{contextual intimacy}: deep, trust-based understanding of how a specific client needs to evolve within its specific ecosystem --- given its politics, its legacy systems, its risk appetite, and its competitive dynamics.

This contextual knowledge has four properties that make it resistant to commoditization. It is too embedded in the client's organizational reality to transfer to a hyperscaler. It requires years of relationship to accumulate, making it unreplicable by a large integrator on a project timeline. It requires judgment across organizational boundaries that no product vendor can automate. And it depends on access to implicit organizational knowledge --- the politics, the trust relationships, the history of failed initiatives, the unspoken constraints --- that does not exist in any training corpus and cannot be reproduced by generative AI.

Contextual intimacy is not a static asset. It compounds: every interaction deepens understanding, every stewardship act increases trust, increased trust increases access, and increased access increases knowledge. This flywheel is the economic logic of stewardship --- and it is precisely the flywheel that extraction destroys.

\subsection{Research Question and Contribution}

This paper addresses the theoretical gap through abductive theory building \citep{DuboisGadde2002}: starting from the empirical anomaly, systematically demonstrating theoretical inadequacy through a structured literature review, and proposing a new construct. Our research question:

\begin{quote}
\emph{How do IT suppliers orchestrate innovation ecosystems without platform ownership, and what capabilities and governance mechanisms enable this role?}
\end{quote}

We introduce \textbf{Ecosystem Stewardship} as a theoretical construct defined by four mechanisms: External Paradox Resolution, Portfolio Transfer, Boundary Management, and Innovation Risk Orchestration. Each mechanism extends an existing theoretical stream by relaxing a specific assumption that blocks explanation of the observed behavior. The construct is developed following \citeauthor{Whetten1989}'s (\citeyear{Whetten1989}) criteria for theoretical contribution and grounded in a systematic literature review of [N] papers across five theoretical streams.

The contribution is threefold. First, we challenge the assumption --- shared across TCE, DC theory, and ecosystem strategy --- that effective orchestration requires either platform control or contractual authority. Second, we extend stewardship theory from intra-firm to inter-organizational relationships, providing a theoretical foundation for the empirically observed but theoretically unexplained behavior of external actors who steward client ecosystems. Third, we provide boundary conditions specifying when ecosystem stewardship creates value (complex ecosystems, capability gaps, trust-based relationships) and when it fails (commoditized services, power asymmetry, transactional relationships).

The paper proceeds as follows. Section II describes our abductive methodology and systematic literature review protocol. Section III synthesizes five literature streams, demonstrating the theoretical inadequacy of each. Section IV presents the Ecosystem Stewardship construct with its four mechanisms and testable propositions. Section V discusses implications, boundary conditions, and limitations. Section VI concludes with an empirical agenda for validation.
